\chapter{Discusión} \label{chp:discussion}

%%%%%%%%%%%%%%%%%%%%%%%%%%%%%%%%%%%%%%%%%%%%%%%%%%%%%%%%%%%%%%%%%%%%%%%%%%%%%%%%
%%%%%%%%%%%%%%%%%%%%%%%%%%%%%%%%%%%%%%%%%%%%%%%%%%%%%%%%%%%%%%%%%%%%%%%%%%%%%%%%

\section{Rendimiento frente a interpretabilidad}

El árbol irrestricto casi memoriza el entrenamiento y produce miles de hojas,
por lo que no es adecuado para explicar decisiones. El árbol final ofrece un
compromiso más útil: conserva un rendimiento similar al modelo con todas las
variables y resulta mucho más fácil de inspeccionar.

La selección de variables favorece una explicación breve, aunque variables
correlacionadas podrían producir subconjuntos diferentes. Además, la selección
no está anidada en una evaluación externa, por lo que el resultado debe
considerarse una estimación de desarrollo.

\section{Limitaciones}

El umbral elegido mejora la detección de dificultades a costa de más falsas
alarmas. Su conveniencia depende del coste relativo de ambos errores y debería
revisarse antes de un uso operativo.

Por último, las reglas describen asociaciones del predictor, no causas del
impago. Esto es especialmente relevante para la asistencia de reembolso, que
puede reflejar información previa sobre la situación del cliente.
